\section{The Evil Demiurge of the Modern World}

In Gnosticism, the Demiurge was the creator of the material world. In opposition to the God, the Demiurge is evil and malevolent, creating a world that hides the true God by entrapping souls in illusion and materiality. The world picture of the modern world is the result of believing the Demiurge.

Hermetic practice must begin with a process of unknowing. That is, one's habitually condition modes of thought must be deconditionalized. Everything assumed or taken for granted must be rethought.

Ideas, worldviews, or ``memes'' enter into the thinking of the world, or noosphere, in three stages:

\begin{enumerate}


\item The first tier of intelligence creates a new system of thought. Such thinkers may be wickedly intelligent. They latch onto a fixed idea and then develop an architectonic system based on it.

\item The second tier absorbs and propagates the ideas from the first tier. These thinkers are in a position to understand the ideas of the first tier. Since ideas are fundamental to them, they are eager to seek out new ideas. Motivated by the need for novelty, intellectual gluttony, and a revolutionary spirit, they draw out the conclusions and diffuse the ideas through publications and educational institutions.

\item The third tier then accepts the ideas uncritically. They have sufficient intelligence to learn the ideas yet lack the depth of insight to reach the roots of the worldview. For this group, certain ideas act as ``status markers'', that is, they are used to separate the ``educated'' from the masses. Articulateness is more important than depth of understanding.
\end{enumerate}


Julius Evola calls the process of unknowing the Trial by Fire\footnote{\url{https://gornahoor.net/?p=1902}}. When done properly, it can seem to lead to the brink of insanity although there are sufficient buffers to avoid going over the edge for a healthy person. I can describe that process at another time, but the first step is to understand the sources of one's thinking.

\paragraph{Founding Fathers}


When I was a young man, I managed to forge a world view founded on Physics, \textbf{Charles Darwin}, \textbf{Sigmund Freud}, and \textbf{Karl Marx}. This would provide a complete explanation of the World: it accounts for matter, life, mind, and society. Around the ages of 10 and 11, I was sent to weekly science classes at the Boston Museum of Science. While waiting to be picked up, I spent time at the library, reading books mostly on astronomy and mathematics, with a particular attachment to the four volume \emph{World of Mathematics} by James Newman. At home, I had a small library of books on astronomy, physics, space travel, rocketry, and dinosaurs, so I was well-disposed to make science the bedrock of my worldview.

Later, at University, I was exposed to the ideas of the New Left. I spent many hours reading turgid Marxian analyses of current events. The more interesting books combined Marx and Freud: those by \textbf{Wilhelm Reich}, \emph{Eros and Civilization} by \textbf{Herbert Marcuse}, and \emph{Life Against Death} by \textbf{Norman O Brown}. That completed the sequence as I read more deeply in psychoanalysis: Freud, Jung, Adler, Fromm, Horney, and so on.

Life on Earth is motivated by Fear, Sex, and Hunger, which maintain us in bondage. This constructed worldview promised liberation from all that. Unlike Buddhism which teaches the desire leads to suffering, the Modern view is the repression of desire leads to suffering. Repression comes from the superego on the individual level and from the patriarchal hierarchy on the social level. So all sexual expression is healthy. I went from high school where, in guidance classes, the issue of premarital sex was still unresolved, to university where the \emph{Harrad Experiment} was a popular book. Oddly enough, the free sex of the Experiment has morphed into contractual sex on campuses where any breach of contract can lead to a Kafkaesque inquisition. Nevertheless, certain practices that were considered immoral, unhealthy, and illegal are now generally accepted as normal. The surprise is the persistence of just a few sexual taboos, but that may yet change. We are left now with greater incidences of incurable venereal diseases, date rape, illegitimacy, abortion, pornography, and so on. Whether all this is a sign of ``liberation'' or an indication of a general malaise is an exercise left to the reader.

In the social realm, Marxism was supposed to lead to abundance, thus solving the problem of hunger (in the larger sense). Since consciousness was formed by the material conditions of life, man's mind would be freed for higher pursuits. The worker would return home for discussions in the High Culture of art, philosophy, and so on. That never happened\footnote{\url{https://gornahoor.net/?p=1582}}. Nowadays, the worker returns home, gets stoned or drunk, and then watches mindless TV shows, sports spectacles, and Internet porn.

The third motivation, Fear, has been intractable because the problem of death has not been solved. \textbf{Ernest Becker} described this in psychoanalytic terms in the \emph{Denial of Death}. The Existentialists brought doubt into the project of liberation via sexual restraint and material abundance. The Transhumanists and Cosmists believe that this problem, too, will be solved. That post is scheduled for the near future.

\paragraph{The Secret Mysterian}


Oddly enough, I also had a Mysterian thread in my thinking life that was nourished at the same time. It can actually be traced back to childhood, but that is another discussion. The brother of a friend used to give out copies of \emph{The Book}, by \textbf{Alan Watts}, which I read. It was unlike anything I had been familiar with, and his other books were my first introduction to the name Rene Guenon. Somehow I balanced this type of neo-Vedanta with an interest in science and psychoanalysis. I believe that my preference for spicy and picante foods prevented me from falling completely into Cogni-ism.

Nevertheless, at the third tier, the four pillars of science, evolution, psychology, and liberalism/socialism, perhaps tinged with some New Age spirituality, even Cogni-ism, supports the worldview of the ``educated'' person today. Thom Hartmann, whose talk show inscrutably airs nightly on Russia Today TV, is one of the more self-aware exemplars of that.

Part 2 will feature the tactics of the Demiurge.


\flrightit{Posted on 2015-10-28 by Cologero}

\begin{center}* * *\end{center}

\begin{footnotesize}
\begin{sffamily}
\texttt{Aleksandar Gueric on 2015-10-29 at 03:32 said:}

``World picture of the modern world is the result of believing in the Demiurge.'' This striked me as very important because I've seen a lot of neo-gnostics in alternative media lately promoting gnostic worldview, this simplified Dualism. It's a Hollywood style of good guys vs bad guys type of thing. I was puzzled by the Demiurge question for a while in the past. Starting with Plato, he obviously didn't talked about the Demiurge that Gnostics later had in mind.

Is it the personification of what Guenon called Countertradition? Personified Maya? Or Samsara? Whatever it is, I understand it's a part of the cycle, of Oneness, becoming to us the unnecessary evil but it's only a shift in balance of rising tides of antitradition during Kali yuga, not some cosmic alien overlord leading the army of Archonts and wanting to overthrow God. Dualism of two principles isn't just simplified version but a fundamentaly wrong version. This may as well be the source of maya.

This Demiurge is unnecessary not just in some abstract cosmic cycle way but for our own realization. It has a practical, personal side. But I wonder what comes after those Evolian trials by fire, suffering and love?

Also, I see the Left as a manifestation of this sinister current. Luciferian, revolutionary, Promethean and eventually self binding. It opposes every social restraint and hierarchical order set by the Right or what they call Patriarchy.

Great article, can't wait for the second part.

\hfill

\texttt{Thomas Blanchard on 2015-10-29 at 11:44 said:}

Powerful writing.

With regard to how the Marxian future was supposed to be, it's interesting to look at science fiction. ``Star Trek: The Next Generation'' is probably the clearest exposition of this vision of the future, where men don't have to work, money is abolished, and all expend their energies and lives ``improving themselves'' and becoming more cultured. It's an attractive vision, if unrealistic. However, I think that for well-oriented person living today, it \_is\_ possible to profit by taking this approach toward life. Much less effort is required today for survival and sustenance -- one can either be seduced and distracted like the masses, or choose to use the unparalleled opportunity for self-development.

Another interesting dichotomy that science fiction has illustrated is that between a future in which man has developed, but dominated and subordinated his technology, his spirituality and heroism acting as primary foundation of civilization, and a future in which man is as crude and vulgar (or perhaps moreso) as he is today, and is completely reliant on his machines and gadgets for his brute power. This dichotomy was depicted most starkly in 1997 videogame, ``StarCraft'', in the contrast between the Protoss (\url{https://www.youtube.com/watch?v=-DJ2r01oOUE}) and Terran (\url{https://www.youtube.com/watch?v=4tdTdnrLVa0}) races.

\hfill

\texttt{klamuse on 2015-10-29 at 14:46 said:}

Abstract beauty in Plato, or the Idea, is not the highest beauty but is in reality only a secondary definition of the most beautiful material object, or supermaterial objects of Godhood at the zenith of evolution. As in Plato, Godhood remains the point of reference for absolute beauty, and for goodness, but it is not the non-material ``beauty'' of Plato. The Great Spiritual Blockade to real beauty and real Godhood is this way unblocked in art, as well as in philosophy and religion, as we understand that we evolve to Godhood in the material and supermaterial world.

Plato says only once in his writings (according to Whitney Oates) that God creates the sacred Ideas, but in the rest of Plato he says life is the creation of the Demiurge, which is right in line with the esoteric understanding in traditional religion which rejects materialism by following the Inward Path and not the Outward Path to Godhood. To them God is not material or supermaterial but non-material---God is a sacred Idea or Word, or a so-called non-material blissful symbolic experience.

Contrary to Plato, life evolves to Godhood in the material and supermaterial world... The big difference from Plato and the traditionalists comes in my defining Godhood as material and supermaterial, as that which we evolve to become in the material world.

This is a different religious understanding from the ones you usually hear, I hope you give it fair attention.

K. L. Anderson

\hfill

\texttt{Traveller on 2015-10-29 at 15:11 said:}

Thanks for your, as usual, excellent insights and observations. The ``demiurge'' element has been a problem floating around in my thoughts for a while now (i.e., is it just conceptual reification based on attitude), so it was interesting to see it come up here.

One comment: Your ``fear, sex and hunger'' are of course included in the first of the three poisons in Buddhism --- attachment (the other two being anger and ignorance), but only to maintain the sense that Buddhism ``teaches desire leads to suffering'' has always seemed incomplete, and a more contextualized sense would actually turn to attachment itself as that which is already part of our existential predicament (with desire only one kind of attachment). Hence, suffering is that which must necessarily be endured (to suffer = to endure) but not be attached to the experience of it. Perhaps this is why, for instance, it is said that bodhicitta should be extended towards the self even before it is extended to the Other.

Also, it's worth pointing out the Buddhist sense in which attachment and karma are connected; that we're only here through our many transmigrations because we haven't yet exhausted all of our karmic streams. However, it does seem that many beings actually generate karma for the sake of experience, just as some children are drawn to things under prohibition, like fire, in order to learn from getting burned. So, apropos of your comments around taboos and disease, as well as various dead ends of intellectual pursuits, the movement of many through various forms of malaise which you cited is an important part of their learning curve... thankfully, some of us no longer need the lessons.

As for the Demiurge, in this light, although it may well be a result of dualistic vision, it perhaps is necessary for some beings, in reasons suggested above, for working through their ``stuff,'' whether as conscious projection or unconscious (karmic) manipulation... who knows? The true scope of the field of interdependence, sometimes called reality, is vast and largely inaccessible, although science likes to think it possesses a high degree of legitimate access, while taking its own limitations to be the furthest extent of intelligence... In the end, perhaps intuition can be valuable for sorting through these kinds of things.

\hfill

\texttt{Bryan M. M. Reynolds on 2015-10-29 at 16:01 said:}

Thanks, Cologero

\hfill

\texttt{Exú on 2015-10-29 at 20:15 said:}

Aleksandar

``But I wonder what comes after those Evolian trials by fire, suffering and love?''

What comes after is the acquire of a self,of a permanent center,that will function in union and harmonie with existence

\hfill

\texttt{yourusernameistaken on 2015-10-30 at 03:30 said:}

This was a powerful post for me. I am undertaking a MsC, and being of a ``spiritual'' bent, I often feel as if higher education is a double-edged sword. One could take it to good place, being motivated in service, or one could use it to secure a good job where one could... as said, work, return home and get drunk and stoned only to watch porn.

\hfill

\texttt{Political Prisoner on 2015-10-30 at 10:01 said:}

Why are there always fluff pieces posted here consisting mostly of petty attacks and a sprinkling of actual information? You've been playing to the crowd cologero. You should have your own talk show, you're a natural.

\hfill

\texttt{Mark Citadel on 2015-10-30 at 13:34 said:}

It is profound in biblical terms that we are asked to choose between God and the world. Some may see this as a horrific choice imposed upon man, but I feel exactly the opposite. It is liberating in the most extreme sense. Serving greater ends puts the world in its true place, below the feet of the Divine Realm. Serve God, and the world is unmasked as the enemy.

I have for well over a year now spoken against the Transhumanist dreams of eliminating death and thus fear. Utopian fantasia at its worst. Why would we not embrace death, unless we were not saved and redeemed? Conquer death? As foolish as thinking you can escape God.

\hfill

\texttt{Tulkas on 2015-10-30 at 16:54 said:}

``I am the source of all spiritual and material worlds. Everything emanates from Me. The wise who know this perfectly engage in My devotional service and worship Me with all their hearts.''

The notion that the material world is the creation of an evil false god is unique to gnosticism. The major world religions teach us that god created the earthly realm. The gnostic mistake is likely what led to the current revolt against nature (human nature, gender, race, heterosexuality etc) and the planned synthetic re-write of nature (genetic modification, geo-engineering, cyborgism). If nature and the human body are the work of an evil entity rather than a loving god, why not improve them?

\hfill

\texttt{Alistair Fraser on 2015-10-30 at 20:15 said:}

Here are my words of personalised meditations on the demiurgical force for this evening, using the terminology of the i-consciousness as the signifier of the individual personas now manifested on this planetarium.

There is a perfectly understandable and simplistic model for a demiurgical force in this realm of the planetarium.

Each i-consciousness has a negated or evolved sporadic or intentionalised direct connection into the psychic realm which is located in parallel with the visible planetarium and does affect directly to indirectly (time proportionate depending on evolved certainty) what shall manifest into the realm of their physical senses.

It can be speculated that there is an overseer of control in this psychic realm but also that there is a law of development in that the personas that take the care to evolve their understanding can also evolve their conscious psychic influence on ``occurrences to be'' It is also possible that just being a good persona in a non-conscious intuited manner can produce similar improvements, although much more ltd in comparison to the ``hero''

Now bearing in mind that one either believes that their life has or has no eternal purpose , this belief also directly affects what ``occurs '' based on the ambiguity or sincerity of that evolved belief in the persona.

Let us then put forth the proposition that the unspoken undefinable overseer has put in place a form of test for these inhabitants of the planetarium which not only spans all passing of time , space and culture, but also spans all levels of development of the psyche.

And this test manifests through the apparent to very real (proportionately believable by the order of ones investment into occurrences to be) by using the analogy of at every moment in ones existence, there is placed before one, a rabbit hole or rabbit holes that offer up excellent rewards both instantly perceivable or some more intricate in their perceived rewards .

Some of these holes are genuine and contain solutions to progress, others are the demiurge imposing its absolute knowledge of you to you , but you can't quite see or even believe in this possibility unless you are a ``hero''

The first sign of the hole is through an apparent spontaneous occurrence which catches the attention and temptates the i-consciousness , once the i-consciousness locks on to the lure, the hole begins to be dug by they , now here is the important part, either on first sniff, first dig, or some ways down that hole, the i-consciousness will by way of its direct connection into the psychic realm receive a warning that something does not feel right , (bearing in mind that the demiurge can also confuse feelings to spoil an immense opportunity that it sees before thee) this is the moment where a lesson is learned or ignored, and so these lures continue, not just for the unaware i-consciousness but equally for the very well developed i-consciousnesses because they can be seen as larger prey for the demiurgical force.

In this manner of ongoing potential traps set in place by the demiurge under the absolute authority of the overseer , a persona is constantly tested in their sincerity of heart and clarification of purpose.

Then as both sincerity of heart and clarification of purpose are both in themselves attenuating conditions ,through both psychic and physical stimuli , the path of life is much more complex and challenging than is recorded .

Let me just project a motion of the hero that has evolved its i-consciousness to a level that demands (occurrences fall to its pleasing) when the timing is correct, the hero can do great works with incredible influence, but within this works also is agitated an excitation of the heroes ego, and if they fail to recognise and defuse this excitation , then the demiurge rubs its hands as the excitation can blind the hero to the subliminised absolute co-ordination of occurrence and the hero dares to think whilst in this excitated propulsion and therein lays heightened danger for ``occurrences'' as the demiurge rubs its hands ... .in excitement , in that very hero , there is the riddle of the thing in and off itself .

Silence is golden at the decisive crossings whispered the grand old lady to thine.

\hfill

\texttt{White Rabbit on 2015-10-31 at 00:24 said:}

Tulkas: that is certainly a unique theory. However, if it's true you will have to explain why modernity is so hedonistic, materialistic, and aspiritual. Anyways, because it does not discriminate, egalitarianism actually extends its love to the whole world, so it hardly opposes the Demiurge.

As for engineering the world, the state of nature is brutal. Should we live like animals in the woods? Why not improve it?

\hfill

\texttt{Logres on 2015-11-11 at 19:31 said:}

``The Transhumanists and Cosmists believe that this problem, too, will be solved. That post is scheduled for the near future... ''

I remember the first time I browsed a copy of The Promethean, in which it was claimed that man would download his consciousness onto computer chips, thus living forever and conquering Death. Someone has been watching too many re-runs of Tron. As an acquaintance of mine remarked, if it were done, the first message we would get back, recycled 30 million times a second, would be ``Turn It OFF''.

\hfill

\texttt{Bryan M. M. Reynolds on 2015-11-12 at 07:26 said:}

@Logres, professor John Gray evidently develops such a thesis about modern gnosticism in his book \_The Soul of the Marionette\_. I haven't read it myself but it sounds interesting.

\hfill

\texttt{Akira on 2015-12-09 at 22:19 said:}

\begin{quote}

although there are sufficient buffers to avoid going over the edge for a healthy person
\end{quote}


This assumes that the products of the modern west are in any way healthy.

Approach the trial by fire with great care, and assume nothing.

Those who attempt it without a strong hermetic or meditative discipline are risking a direct experience with the negative principle that they literally may not survive.


\hfill

\end{sffamily}
\end{footnotesize}

